\uuid{UoKo}
\titre{Espaces vectoriels}
\auteur{hanani ; mzari}
\organisation{exo7}
\licence{CC BY-NC-SA 4.0}
\datecreate{2020-01-20}
\niveau{L1}
\module{Algèbre}
\chapitre{Espace vectoriel}
\sousChapitre{Définition, sous-espace}
\difficulte{2}
\theme{Espaces vectoriels}
\exo7qcmid{lille-2.espaces-vectoriels.006}

\contenu{
\begin{question}
\qclassification{106.01}
\qtags{motcle=Sous-espace vectoriel de $\Rr^n$}
Soit $E=\{(x,y,z) \in \Rr^3 \,  ; \; (x+y)(x+z)=0\}$, muni des opérations usuelles. Quelles sont les assertions vraies ?
\begin{answers}
\bad{$E=\{(x,y,z) \in \Rr^3 \, ; \; x+y=x+z=0\}$.}
\bad{$E=\{(x,y,z) \in \Rr^3 \, ; \;  x+y=0\} \cap   \{(x,y,z) \in \Rr^3 ; \; x+z=0\}$.}
\good{$E=\{(x,y,z) \in \Rr^3 \, ; \;  x+y=0\} \cup   \{(x,y,z) \in \Rr^3 ; \; x+z=0\}$.}
\good{$E$ n'est pas un espace vectoriel, car $E$ n'est pas stable par addition.}
\end{answers}
\begin{explanations}
$E= \{(x,y,z) \in \Rr^3 \,  ;\;  x+y=0 \; \mbox{ou} \;  x+z=0  \} =\{(x,y,z) \in \Rr^3 \,  ;\;  x+y=0\} \cup
\{(x,y,z) \in \Rr^3 \,  ; \;  x+z=0\}$.
$E$ n'est pas un espace vectoriel, car $E$ n'est pas stable par addition : $(1,-1,0), (1,0,-1) \in E$, mais,
$(2,-1,-1) \notin E$.
\end{explanations}
\end{question}
}
